%% toc: P2 16 16.5,16.6,16.7
\section{Proofs for the elimination of a defined relation symbol}
\label{sec:annex-p02-w02b-p02-018-defined-relation-elimination}

The standing data are those of \cref{def:p02-w02b-p02-018-standing}.  Two
facts about assignment updates are used repeatedly; both follow directly
from the definition of \(s[x\mapsto a]\) in \cref{ch:fol-semantics}.

\emph{Fact U1.}  If \(u\ne v\), then
\(s[u\mapsto a][v\mapsto b]=s[v\mapsto b][u\mapsto a]\): both assignments
send \(u\) to \(a\), \(v\) to \(b\), and every other variable \(w\) to
\(s(w)\).

\emph{Fact U2.}  If \(v_1,\ldots,v_k\) are pairwise distinct, then
\(s[v_1\mapsto b_1]\cdots[v_k\mapsto b_k]\) sends \(v_i\) to \(b_i\) for
each \(i\) and agrees with \(s\) on every other variable.  By induction on
\(k\): the case \(k=0\) is trivial; the last update sets \(v_k\) to
\(b_k\) and leaves every other variable at its previous value, which by
the induction hypothesis is \(b_i\) at \(v_i\) for \(i<k\), since
\(v_i\ne v_k\), and \(s(w)\) at every other \(w\).

\begin{annexitem}{annex:p02-w02b-p02-018-instance}{thm:p02-w02b-p02-018-instance}{Instance lemma}
\begin{proof}
Write \(V\), \(z_1,\ldots,z_n\), \(\chi_0,\ldots,\chi_n\) and
\(\theta_0,\ldots,\theta_n\) as in
\cref{def:p02-w02b-p02-018-instance}.  By the choice of the \(z_k\):
\begin{equation*}
 \begin{aligned}
  &z_1,\ldots,z_n\text{ are pairwise distinct};\\
  &z_k\notin\{x_1,\ldots,x_n\}\text{ and }
   z_k\notin\operatorname{Var}(t_j)\text{ for all }j,k.
 \end{aligned}
 \tag{\(\dagger\)}
\end{equation*}

\emph{Clause 1, membership in \(\mathsf{Fm}(\Lang)\).}  First, if
\(r,t\in\mathsf{Term}(\Lang)\), then \(r[t/x]\in\mathsf{Term}(\Lang)\): by
induction on \(r\), a variable becomes \(t\) or stays a variable, and
\(f(r_1,\ldots,r_m)[t/x]=f(r_1[t/x],\ldots,r_m[t/x])\) applies the function
symbol \(f\) of \(\Lang\) to \(\Lang\)-terms by the induction hypothesis.
Next, if \(\chi\in\mathsf{Fm}(\Lang)\) and \(t\in\mathsf{Term}(\Lang)\), then
\(\chi[t/x]\in\mathsf{Fm}(\Lang)\), by strong induction on \(c(\chi)\).  An
atomic formula keeps its equality sign or its relation symbol of \(\Lang\)
and receives \(\Lang\)-terms.  The clauses for \(\bottom\),
\(\neg\), \(\land\), \(\lor\) and \(\to\) apply the same constructor to
results that lie in \(\mathsf{Fm}(\Lang)\) by the induction hypothesis.  For
\(\chi=Qy\,\eta\), the first substitution clause returns \(\chi\); the
second returns \(Qy\,(\eta[t/x])\) with \(c(\eta)<c(\chi)\); the third
returns \(Qz\,((\eta[z/y])[t/x])\), where \(\eta[z/y]\in\mathsf{Fm}(\Lang)\)
by the induction hypothesis at \(c(\eta)\), and
\(c(\eta[z/y])=c(\eta)<c(\chi)\) by \cref{lem:substitution-complexity}, so
the induction hypothesis applies once more.  Since \(\varphi\in\mathsf{Fm}(\Lang)\)
and every \(z_k\) and \(t_k\) is an \(\Lang\)-term, each \(\chi_k\) and each
\(\theta_k\) lies in \(\mathsf{Fm}(\Lang)\) by induction on \(k\).

\emph{Clause 1, free variables.}  We show by induction on \(k\) that
\[
 \FV(\chi_k)\subseteq\{z_1,\ldots,z_k\}\cup\{x_{k+1},\ldots,x_n\}.
\]
For \(k=0\) this is the hypothesis on \(\varphi\).  For \(k\geq1\),
\cref{lem:substitution-fv}(1) gives
\(\FV(\chi_k)\subseteq(\FV(\chi_{k-1})\setminus\{x_k\})\cup\{z_k\}\).
By the induction hypothesis and \((\dagger)\), removing \(x_k\) from
\(\{z_1,\ldots,z_{k-1}\}\cup\{x_k,\ldots,x_n\}\) leaves a subset of
\(\{z_1,\ldots,z_{k-1}\}\cup\{x_{k+1},\ldots,x_n\}\), because \(x_k\) is not
among the \(z_j\) and differs from \(x_{k+1},\ldots,x_n\).  Adding \(z_k\)
gives the claim.  Thus \(\FV(\theta_0)\subseteq\{z_1,\ldots,z_n\}\).  Next,
by induction on \(k\),
\[
 \FV(\theta_k)\subseteq\{z_{k+1},\ldots,z_n\}
 \cup\operatorname{Var}(t_1)\cup\cdots\cup\operatorname{Var}(t_k).
\]
For \(k\geq1\), \cref{lem:substitution-fv}(1) gives
\(\FV(\theta_k)\subseteq(\FV(\theta_{k-1})\setminus\{z_k\})
\cup\operatorname{Var}(t_k)\).  Removing \(z_k\) from the bound for
\(\FV(\theta_{k-1})\) leaves a subset of
\(\{z_{k+1},\ldots,z_n\}\cup\operatorname{Var}(t_1)\cup\cdots
\cup\operatorname{Var}(t_{k-1})\), since the \(z_j\) are distinct.  Adding
\(\operatorname{Var}(t_k)\) gives the claim, and \(k=n\) gives the stated
bound for \(\varphi(t_1,\ldots,t_n)=\theta_n\).

\emph{Clause 2.}  Fix \(\mathcal N\), \(s\) and
\(a_i=\sem{t_i}^{\mathcal N}_s\).  All formulas involved are
\(\Lang\)-formulas, hence formulas of the signature of \(\mathcal N\), so
\cref{lem:fol-coincidence}, \cref{thm:fol-substitution} and term coincidence
apply to them in \(\mathcal N\).

\emph{Step A, insertion stage.}  For \(k=0,\ldots,n\) let
\(s_k\defeq s[z_n\mapsto a_n][z_{n-1}\mapsto a_{n-1}]\cdots
[z_{k+1}\mapsto a_{k+1}]\), so that \(s_n=s\).  We show by downward
induction on \(k\) that
\(\mathcal N,s\entails\theta_n\) iff \(\mathcal N,s_k\entails\theta_k\).
The case \(k=n\) is trivial.  Suppose the equivalence holds for some
\(k\geq1\).  Since \(\theta_k=\theta_{k-1}[t_k/z_k]\),
\cref{thm:fol-substitution} gives
\(\mathcal N,s_k\entails\theta_k\) iff
\(\mathcal N,s_k[z_k\mapsto\sem{t_k}^{\mathcal N}_{s_k}]\entails
\theta_{k-1}\).  By Fact U2, \(s_k\) differs from \(s\) at most at
\(z_{k+1},\ldots,z_n\), none of which lies in \(\operatorname{Var}(t_k)\) by
\((\dagger)\); term coincidence gives
\(\sem{t_k}^{\mathcal N}_{s_k}=\sem{t_k}^{\mathcal N}_s=a_k\).  Hence the
updated assignment is \(s_{k-1}\), and the equivalence holds for \(k-1\).
At \(k=0\):
\[
 \mathcal N,s\entails\varphi(t_1,\ldots,t_n)
 \quad\Longleftrightarrow\quad
 \mathcal N,s_0\entails\chi_n .
\]

\emph{Step B, renaming stage.}  We show by induction on \(k\) that for
every assignment \(u\) in \(\mathcal N\),
\[
 \mathcal N,u\entails\chi_k
 \quad\Longleftrightarrow\quad
 \mathcal N,u[x_1\mapsto u(z_1)]\cdots[x_k\mapsto u(z_k)]\entails\varphi .
\]
The case \(k=0\) is trivial.  Let \(k\geq1\) and let \(u\) be given.  Since
\(\chi_k=\chi_{k-1}[z_k/x_k]\) and \(\sem{z_k}^{\mathcal N}_u=u(z_k)\),
\cref{thm:fol-substitution} gives \(\mathcal N,u\entails\chi_k\) iff
\(\mathcal N,u'\entails\chi_{k-1}\), where
\(u'\defeq u[x_k\mapsto u(z_k)]\).  By the induction hypothesis at \(u'\),
this holds iff
\(\mathcal N,u'[x_1\mapsto u'(z_1)]\cdots[x_{k-1}\mapsto u'(z_{k-1})]
\entails\varphi\).  For \(j<k\), \(z_j\ne x_k\) by \((\dagger)\), so
\(u'(z_j)=u(z_j)\).  The variables \(x_1,\ldots,x_k\) are pairwise distinct,
so repeated use of Fact U1 moves the update at \(x_k\) past the updates at
\(x_1,\ldots,x_{k-1}\):
\[
 \begin{aligned}
 &u[x_k\mapsto u(z_k)][x_1\mapsto u(z_1)]\cdots[x_{k-1}\mapsto u(z_{k-1})]\\
 &\quad=u[x_1\mapsto u(z_1)]\cdots[x_k\mapsto u(z_k)].
 \end{aligned}
\]
This proves the claim for \(k\).

\emph{Combination.}  By Fact U2 and \((\dagger)\), \(s_0(z_i)=a_i\) for
every \(i\).  Step B with \(k=n\) and \(u=s_0\) gives
\(\mathcal N,s_0\entails\chi_n\) iff
\(\mathcal N,s_0[\bar x\mapsto\bar a]\entails\varphi\).  By Fact U2 the
assignments \(s_0[\bar x\mapsto\bar a]\) and \(s[\bar x\mapsto\bar a]\) both
send each \(x_i\) to \(a_i\), so they agree on
\(\FV(\varphi)\subseteq\{x_1,\ldots,x_n\}\), and
\cref{lem:fol-coincidence} gives
\(\mathcal N,s_0[\bar x\mapsto\bar a]\entails\varphi\) iff
\(\mathcal N,s[\bar x\mapsto\bar a]\entails\varphi\).  Chaining with Step A
proves clause 2.

\emph{Clause 3.}  If \(\bar a\in R^{\mathcal M}_{\varphi}\), then
\(\mathcal M,s[\bar x\mapsto\bar a]\entails\varphi\) by the definition,
which quantifies over every assignment.  Conversely, suppose
\(\mathcal M,s[\bar x\mapsto\bar a]\entails\varphi\) and let \(s''\) be any
assignment.  By Fact U2, \(s''[\bar x\mapsto\bar a]\) and
\(s[\bar x\mapsto\bar a]\) both send each \(x_i\) to \(a_i\), so they agree
on \(\FV(\varphi)\); \cref{lem:fol-coincidence} gives
\(\mathcal M,s''[\bar x\mapsto\bar a]\entails\varphi\).  As \(s''\) was
arbitrary, \(\bar a\in R^{\mathcal M}_{\varphi}\).
\end{proof}
\end{annexitem}

\begin{annexitem}{annex:p02-w02b-p02-018-translation-syntax}{thm:p02-w02b-p02-018-translation-syntax}{The translation lands in the old signature}
\begin{proof}
Both clauses are proved by strong induction on \(c(\psi)\), following the
clauses of \cref{def:p02-w02b-p02-018-translation}.  Since \(\Lang\) and
\(\Lang_R\) have the same function symbols, every term occurring in
\(\psi\) is an \(\Lang\)-term.

\emph{Clause 1.}  If \(\psi\) is \(r=t\) or \(S(t_1,\ldots,t_m)\) with
\(S\in\Lang\), then \(\psi^{*}=\psi\) is an atomic \(\Lang\)-formula, and
the free variables are unchanged.  If \(\psi=R(t_1,\ldots,t_n)\), then
\(\psi^{*}=\varphi(t_1,\ldots,t_n)\), which lies in \(\mathsf{Fm}(\Lang)\) and
has free variables among
\(\operatorname{Var}(t_1)\cup\cdots\cup\operatorname{Var}(t_n)=\FV(\psi)\) by
\cref{thm:p02-w02b-p02-018-instance}(1).  The case \(\bottom\) is
immediate.  For \(\neg\chi\) and \(\chi\circ\eta\), the induction
hypotheses give \(\Lang\)-formulas \(\chi^{*},\eta^{*}\) with
\(\FV(\chi^{*})\subseteq\FV(\chi)\) and \(\FV(\eta^{*})\subseteq\FV(\eta)\).
Then \(\neg\chi^{*}\) and \((\chi^{*}\circ\eta^{*})\) are
\(\Lang\)-formulas, \(\FV(\neg\chi^{*})=\FV(\chi^{*})\subseteq\FV(\neg\chi)\),
and \(\FV(\chi^{*}\circ\eta^{*})=\FV(\chi^{*})\cup\FV(\eta^{*})
\subseteq\FV(\chi)\cup\FV(\eta)=\FV(\chi\circ\eta)\).
For \(Qy\,\chi\), the induction hypothesis gives
\(\chi^{*}\in\mathsf{Fm}(\Lang)\), hence \(Qy\,\chi^{*}\in\mathsf{Fm}(\Lang)\),
and \(\FV(Qy\,\chi^{*})=\FV(\chi^{*})\setminus\{y\}\subseteq
\FV(\chi)\setminus\{y\}=\FV(Qy\,\chi)\).  If \(\psi\) is a sentence, then
\(\FV(\psi^{*})\subseteq\varnothing\).

\emph{Clause 2.}  If \(\psi\in\mathsf{Fm}(\Lang)\), no atomic subformula of
\(\psi\) has the symbol \(R\), because \(R\notin\Lang\).  The atomic clauses
that apply are therefore the identity clauses, \(\bottom^{*}=\bottom\), and
the remaining clauses rebuild \(\psi\) from the translations of its
immediate subformulas, which are themselves \(\Lang\)-formulas and are
fixed by the induction hypothesis.
\end{proof}
\end{annexitem}

\begin{annexitem}{annex:p02-w02b-p02-018-satisfaction-translation}{thm:p02-w02b-p02-018-satisfaction-translation}{Satisfaction translation}
\begin{proof}
Fix the \(\Lang\)-structure \(\mathcal M\).  We prove, by strong induction
on \(c(\psi)\), that for every \(\psi\in\mathsf{Fm}(\Lang_R)\) the
equivalence holds for \emph{every} assignment \(s\); the quantifier case
needs the hypothesis at updated assignments.  Recall that
\(\mathcal M_{\varphi}\mathbin{\upharpoonright}\Lang=\mathcal M\) and that
the two structures have the same domain.

\emph{Equality and old relation symbols.}  For \(\psi\) of the form
\(r=t\) or \(S(t_1,\ldots,t_m)\) with \(S\in\Lang\), \(\psi\) is an
\(\Lang\)-formula and \(\psi^{*}=\psi\); invariance under reduct
(\cref{lem:reduct-invariance}), applied to the \(\Lang_R\)-structure
\(\mathcal M_{\varphi}\) with reduct \(\mathcal M\), gives
\(\mathcal M_{\varphi},s\entails\psi\) iff \(\mathcal M,s\entails\psi\).

\emph{The new symbol.}  Let \(\psi=R(t_1,\ldots,t_n)\) and put
\(a_i\defeq\sem{t_i}^{\mathcal M}_s\).  By \cref{lem:reduct-invariance},
\(\sem{t_i}^{\mathcal M_{\varphi}}_s=a_i\).  The atomic satisfaction clause
and \cref{def:p02-w02b-p02-018-defined-relation} give
\(\mathcal M_{\varphi},s\entails\psi\) iff
\(\bar a\in R^{\mathcal M}_{\varphi}\).  By
\cref{thm:p02-w02b-p02-018-instance}(3) this holds iff
\(\mathcal M,s[\bar x\mapsto\bar a]\entails\varphi\), and by
\cref{thm:p02-w02b-p02-018-instance}(2), with \(\mathcal N=\mathcal M\),
iff \(\mathcal M,s\entails\varphi(t_1,\ldots,t_n)\), that is,
\(\mathcal M,s\entails\psi^{*}\).

\emph{Falsum and connectives.}  Neither structure satisfies \(\bottom\) under
any assignment, and \(\bottom^{*}=\bottom\).  For \(\neg\chi\),
\(\mathcal M_{\varphi},s\entails\neg\chi\) iff
\(\mathcal M_{\varphi},s\not\entails\chi\) iff, by the induction hypothesis
for \(\chi\) at \(s\), \(\mathcal M,s\not\entails\chi^{*}\) iff
\(\mathcal M,s\entails(\neg\chi)^{*}\).

For the binary connectives let \(\psi=(\chi\circ\eta)\) with
\(\circ\in\{\land,\lor,\to\}\).  Then \(c(\chi)<c(\psi)\) and
\(c(\eta)<c(\psi)\), so the induction hypothesis at \(s\) gives
\begin{equation*}
 \mathcal M_{\varphi},s\entails\chi\iff\mathcal M,s\entails\chi^{*},
 \qquad
 \mathcal M_{\varphi},s\entails\eta\iff\mathcal M,s\entails\eta^{*}.
 \tag{IH}
\end{equation*}
\begin{enumerate}
 \item \(\circ={\land}\).  \(\mathcal M_{\varphi},s\entails(\chi\land\eta)\)
       iff \(\mathcal M_{\varphi},s\entails\chi\) and
       \(\mathcal M_{\varphi},s\entails\eta\), by the clause for \(\land\);
       iff \(\mathcal M,s\entails\chi^{*}\) and
       \(\mathcal M,s\entails\eta^{*}\), by (IH); iff
       \(\mathcal M,s\entails(\chi^{*}\land\eta^{*})\), by the clause for
       \(\land\) in \(\mathcal M\); and
       \((\chi^{*}\land\eta^{*})=(\chi\land\eta)^{*}\).
 \item \(\circ={\lor}\).  \(\mathcal M_{\varphi},s\entails(\chi\lor\eta)\)
       iff \(\mathcal M_{\varphi},s\entails\chi\) or
       \(\mathcal M_{\varphi},s\entails\eta\); iff
       \(\mathcal M,s\entails\chi^{*}\) or \(\mathcal M,s\entails\eta^{*}\), by
       (IH); iff \(\mathcal M,s\entails(\chi^{*}\lor\eta^{*})\); and
       \((\chi^{*}\lor\eta^{*})=(\chi\lor\eta)^{*}\).
 \item \(\circ={\to}\).  \(\mathcal M_{\varphi},s\entails(\chi\to\eta)\)
       iff \(\mathcal M_{\varphi},s\not\entails\chi\) or
       \(\mathcal M_{\varphi},s\entails\eta\).  By (IH),
       \(\mathcal M_{\varphi},s\not\entails\chi\) iff
       \(\mathcal M,s\not\entails\chi^{*}\), and
       \(\mathcal M_{\varphi},s\entails\eta\) iff
       \(\mathcal M,s\entails\eta^{*}\); so the condition holds iff
       \(\mathcal M,s\not\entails\chi^{*}\) or \(\mathcal M,s\entails\eta^{*}\),
       that is, iff \(\mathcal M,s\entails(\chi^{*}\to\eta^{*})\), and
       \((\chi^{*}\to\eta^{*})=(\chi\to\eta)^{*}\).
\end{enumerate}

\emph{Quantifiers.}  Let \(\psi=\forall y\,\chi\).  Then
\(\mathcal M_{\varphi},s\entails\psi\) iff
\(\mathcal M_{\varphi},s[y\mapsto a]\entails\chi\) for every
\(a\in\dom\mathcal M\) (the domains coincide).  Since \(c(\chi)<c(\psi)\),
the induction hypothesis at each assignment \(s[y\mapsto a]\) turns this
into: \(\mathcal M,s[y\mapsto a]\entails\chi^{*}\) for every
\(a\in\dom\mathcal M\), which is \(\mathcal M,s\entails\forall y\,\chi^{*}\),
and \(\forall y\,\chi^{*}=\psi^{*}\).

Let \(\psi=\exists y\,\chi\).  By the clause for \(\exists\),
\(\mathcal M_{\varphi},s\entails\psi\) iff there is
\(a\in\dom\mathcal M_{\varphi}=\dom\mathcal M\) with
\(\mathcal M_{\varphi},s[y\mapsto a]\entails\chi\).  For each such \(a\),
the induction hypothesis for \(\chi\), with \(c(\chi)<c(\psi)\), at the
assignment \(s[y\mapsto a]\) gives
\(\mathcal M_{\varphi},s[y\mapsto a]\entails\chi\) iff
\(\mathcal M,s[y\mapsto a]\entails\chi^{*}\).  So
\(\mathcal M_{\varphi},s\entails\psi\) iff there is \(a\in\dom\mathcal M\)
with \(\mathcal M,s[y\mapsto a]\entails\chi^{*}\), which by the clause for
\(\exists\) in \(\mathcal M\) is \(\mathcal M,s\entails\exists y\,\chi^{*}\),
and \(\exists y\,\chi^{*}=(\exists y\,\chi)^{*}\).

\emph{The model-of-\(T\) clause.}  Let \(\mathcal M\) be an
\(\Lang\)-model of \(T\).  By
\cref{thm:p02-w02b-p02-018-consequences}(1), whose proof in
\cref{annex:p02-w02b-p02-018-consequences} does not use the present clause,
\(\mathcal M_{\varphi}\) satisfies \(\delta_R\) and is the only
\(\Lang_R\)-expansion of \(\mathcal M\) that does.  Every \(\sigma\in T\) is
an \(\Lang\)-sentence true in \(\mathcal M\), hence true in
\(\mathcal M_{\varphi}\) by \cref{lem:reduct-invariance}; so
\(\mathcal M_{\varphi}\) is a model of \(T\cup\{\delta_R\}\).  The
equivalence for every \(s\) and \(\psi\) is the case of the main statement
for this \(\mathcal M\).
\end{proof}
\end{annexitem}

\begin{annexitem}{annex:p02-w02b-p02-018-consequences}{thm:p02-w02b-p02-018-consequences}{Consequence and derivability forms}
\begin{proof}
\emph{Clause 1.}  First, \(\delta_R\) is a sentence:
\(\FV(R(x_1,\ldots,x_n))=\{x_1,\ldots,x_n\}\) and
\(\FV(\varphi)\subseteq\{x_1,\ldots,x_n\}\), the defined biconditional
\(((\alpha\to\beta)\land(\beta\to\alpha))\) has the free variables of
\(\alpha\) and \(\beta\), and the \(n\) universal quantifiers remove
\(x_1,\ldots,x_n\).  So its truth in a structure does not depend on the
assignment.

Let \(\mathcal N\) be an \(\Lang_R\)-structure, let
\(\mathcal M\defeq\mathcal N\mathbin{\upharpoonright}\Lang\), and fix an
assignment \(s\).  Applying the satisfaction clause for \(\forall\) once for
each of \(x_1,\ldots,x_n\), \(\mathcal N\entails\delta_R\) iff for every
\(\bar a\in(\dom\mathcal N)^n\),
\(\mathcal N,s[\bar x\mapsto\bar a]\entails
R(x_1,\ldots,x_n)\leftrightarrow\varphi\).  By the clauses for \(\land\) and
\(\to\), a structure satisfies \(((\alpha\to\beta)\land(\beta\to\alpha))\)
under an assignment iff it satisfies both \(\alpha\) and \(\beta\) or
neither.  Fix \(\bar a\) and write \(s'=s[\bar x\mapsto\bar a]\).  By Fact
U2, \(s'(x_i)=a_i\), so the atomic clause gives
\(\mathcal N,s'\entails R(x_1,\ldots,x_n)\) iff \(\bar a\in R^{\mathcal N}\).
By \cref{lem:reduct-invariance} and
\cref{thm:p02-w02b-p02-018-instance}(3),
\(\mathcal N,s'\entails\varphi\) iff \(\mathcal M,s'\entails\varphi\) iff
\(\bar a\in R^{\mathcal M}_{\varphi}\).  Hence \(\mathcal N\entails\delta_R\)
iff, for every \(\bar a\in(\dom\mathcal N)^n\),
\(\bar a\in R^{\mathcal N}\Leftrightarrow\bar a\in R^{\mathcal M}_{\varphi}\);
as both relations are subsets of \((\dom\mathcal N)^n\), this says
\(R^{\mathcal N}=R^{\mathcal M}_{\varphi}\).  The structures \(\mathcal N\)
and \(\mathcal M_{\varphi}\) are both \(\Lang_R\)-expansions of
\(\mathcal M\), so they have the same domain and agree on \(\Lang\); they
are equal iff they agree on \(R\), that is, iff
\(R^{\mathcal N}=R^{\mathcal M}_{\varphi}\).

For the last sentence of clause 1, let \(\mathcal M\) be an
\(\Lang\)-structure.  The structure \(\mathcal M_{\varphi}\) has reduct
\(\mathcal M\) and \(R^{\mathcal M_{\varphi}}=R^{\mathcal M}_{\varphi}\), so
it satisfies \(\delta_R\) by what was just shown.  If \(\mathcal N\) is any
\(\Lang_R\)-expansion of \(\mathcal M\) satisfying \(\delta_R\), then
\(\mathcal N\mathbin{\upharpoonright}\Lang=\mathcal M\) and
\(\mathcal N=\mathcal M_{\varphi}\).

\emph{Clause 2.}  Let \(\mathcal N\) be an \(\Lang_R\)-structure and \(s\)
an assignment with \(\mathcal N,s\entails\delta_R\), and put
\(\mathcal M=\mathcal N\mathbin{\upharpoonright}\Lang\).  By clause 1,
\(\mathcal N=\mathcal M_{\varphi}\).  By
\cref{thm:p02-w02b-p02-018-satisfaction-translation},
\(\mathcal N,s\entails\psi\) iff \(\mathcal M,s\entails\psi^{*}\).  Since
\(\psi^{*}\in\mathsf{Fm}(\Lang)\)
(\cref{thm:p02-w02b-p02-018-translation-syntax}),
\cref{lem:reduct-invariance} gives \(\mathcal M,s\entails\psi^{*}\) iff
\(\mathcal N,s\entails\psi^{*}\).  So \(\mathcal N,s\) satisfies both
\(\psi\) and \(\psi^{*}\) or neither, which by the computation in clause 1
is \(\mathcal N,s\entails\psi\leftrightarrow\psi^{*}\).  This proves
\(\{\delta_R\}\entails\psi\leftrightarrow\psi^{*}\).  Any structure and
assignment satisfying every member of \(T\cup\{\delta_R\}\) satisfy
\(\delta_R\), so \(T\cup\{\delta_R\}\entails\psi\leftrightarrow\psi^{*}\).

\emph{Clause 3.}  Suppose \(T\cup\{\delta_R\}\cup\Gamma\entails\psi\).  Let
\(\mathcal M\) be an \(\Lang\)-structure and \(s\) an assignment satisfying
every member of \(T\cup\Gamma^{*}\).  By clause 1,
\(\mathcal M_{\varphi}\entails\delta_R\); by \cref{lem:reduct-invariance},
\(\mathcal M_{\varphi},s\) satisfies every member of \(T\); and for each
\(\gamma\in\Gamma\), \(\mathcal M,s\entails\gamma^{*}\) gives
\(\mathcal M_{\varphi},s\entails\gamma\) by
\cref{thm:p02-w02b-p02-018-satisfaction-translation}.  Hence
\(\mathcal M_{\varphi},s\entails\psi\), and
\cref{thm:p02-w02b-p02-018-satisfaction-translation} gives
\(\mathcal M,s\entails\psi^{*}\).  Thus \(T\cup\Gamma^{*}\entails\psi^{*}\);
by \cref{thm:p02-w02b-p02-018-translation-syntax}(1) every member of
\(T\cup\Gamma^{*}\cup\{\psi^{*}\}\) is an \(\Lang\)-formula, so this
consequence is one between \(\Lang\)-formulas.

Conversely, suppose \(T\cup\Gamma^{*}\entails\psi^{*}\).  Let \(\mathcal N\)
be an \(\Lang_R\)-structure and \(s\) an assignment satisfying every member
of \(T\cup\{\delta_R\}\cup\Gamma\), and put
\(\mathcal M=\mathcal N\mathbin{\upharpoonright}\Lang\).  By clause 1,
\(\mathcal N=\mathcal M_{\varphi}\).  By \cref{lem:reduct-invariance},
\(\mathcal M,s\) satisfies every member of \(T\); by
\cref{thm:p02-w02b-p02-018-satisfaction-translation}, it satisfies
\(\gamma^{*}\) for every \(\gamma\in\Gamma\).  So
\(\mathcal M,s\entails\psi^{*}\), and
\cref{thm:p02-w02b-p02-018-satisfaction-translation} gives
\(\mathcal N,s\entails\psi\).  Thus
\(T\cup\{\delta_R\}\cup\Gamma\entails\psi\).  The case
\(\Gamma=\varnothing\) is the second displayed equivalence.  If moreover
\(\psi\in\mathsf{Sent}(\Lang)\), then \(\psi^{*}=\psi\) by
\cref{thm:p02-w02b-p02-018-translation-syntax}(2), and the equivalence
reads \(T\cup\{\delta_R\}\entails\psi\) iff \(T\entails\psi\).

\emph{Clause 4.}  Assume the Axiom of Choice.  By clause 2,
\(T\cup\{\delta_R\}\entails\psi\leftrightarrow\psi^{*}\) in the signature
\(\Lang_R\); \cref{thm:fol-completeness}, applied to the signature
\(\Lang_R\) and the formula \(((\psi\to\psi^{*})\land(\psi^{*}\to\psi))\),
gives \(T\cup\{\delta_R\}\derives_{\mathsf{ND}_{\mathsf{FOL}}}
\psi\leftrightarrow\psi^{*}\).  If
\(T\cup\{\delta_R\}\cup\Gamma\derives_{\mathsf{ND}_{\mathsf{FOL}}}\psi\),
then \(T\cup\{\delta_R\}\cup\Gamma\entails\psi\) by
\cref{thm:fol-soundness}, so \(T\cup\Gamma^{*}\entails\psi^{*}\) by
clause 3, and \cref{thm:fol-completeness} in the signature \(\Lang\) gives
\(T\cup\Gamma^{*}\derives_{\mathsf{ND}_{\mathsf{FOL}}}\psi^{*}\).  If
\(T\cup\Gamma^{*}\derives_{\mathsf{ND}_{\mathsf{FOL}}}\psi^{*}\), then
\(T\cup\Gamma^{*}\entails\psi^{*}\) by \cref{thm:fol-soundness}, so
\(T\cup\{\delta_R\}\cup\Gamma\entails\psi\) by clause 3, and
\cref{thm:fol-completeness} in the signature \(\Lang_R\) gives
\(T\cup\{\delta_R\}\cup\Gamma\derives_{\mathsf{ND}_{\mathsf{FOL}}}\psi\).
\end{proof}
\end{annexitem}

\section{Proofs for defined function symbols}
\label{sec:annex-p02-w02b-p02-018-function-symbols}

The standing data are those of \cref{def:p02-w02b-p02-018-function-standing}.
As explained after that definition, \cref{def:p02-w02b-p02-018-instance} and
\cref{thm:p02-w02b-p02-018-instance} are applied to \(\varphi\) with the list
\(x_1,\ldots,x_m,y\); the proof of
\cref{annex:p02-w02b-p02-018-instance} uses only that these variables are
pairwise distinct and contain \(\FV(\varphi)\).  Facts U1 and U2 of
\cref{sec:annex-p02-w02b-p02-018-defined-relation-elimination} are used
again, together with the following one.

\emph{Fact U3.}  For variables \(u_1,\ldots,u_j\) and formulas
\(\alpha_1,\ldots,\alpha_j,\beta\), a structure \(\mathcal M\) and an
assignment \(s\) satisfy
\(\exists u_1\cdots\exists u_j\,(\alpha_1\land\cdots\land\alpha_j\land\beta)\),
with the right-associated conjunction of
\cref{def:p02-w02b-p02-018-graph-formula}, iff there are
\(c_1,\ldots,c_j\in\dom\mathcal M\) such that, writing
\(s_c\defeq s[u_1\mapsto c_1]\cdots[u_j\mapsto c_j]\),
\(\mathcal M,s_c\entails\alpha_i\) for every \(i\) and
\(\mathcal M,s_c\entails\beta\).  First, by induction on \(j\) and the
clause for \(\land\), an assignment satisfies the right-associated
conjunction iff it satisfies each \(\alpha_i\) and \(\beta\): for \(j=0\) the
conjunction is \(\beta\), and for \(j\geq1\) it is
\((\alpha_1\land\gamma)\) with \(\gamma\) the conjunction of
\(\alpha_2,\ldots,\alpha_j,\beta\).  Second, by induction on \(j\) and the
clause for \(\exists\): \(\mathcal M,s\entails\exists u_1\,\eta\) iff
\(\mathcal M,s[u_1\mapsto c_1]\entails\eta\) for some \(c_1\), and the
induction hypothesis applied to the remaining prefix \(\exists u_2\cdots\exists u_j\)
at the assignment \(s[u_1\mapsto c_1]\) supplies \(c_2,\ldots,c_j\).

For a sentence \(\sigma\), \(\mathcal M\entails\sigma\) means
\(\mathcal M,s\entails\sigma\) for some, equivalently every, assignment
\(s\), by \cref{lem:fol-coincidence}.

\begin{annexitem}{annex:p02-w02b-p02-018-unique-value}{thm:p02-w02b-p02-018-unique-value}{The unique-value sentences}
\begin{proof}
\emph{Sentences.}  The \(m+1\) quantifiers of \(\operatorname{Tot}_{\varphi}\)
remove \(x_1,\ldots,x_m,y\supseteq\FV(\varphi)\).  By
\cref{lem:substitution-fv}(1),
\(\FV(\varphi[y'/y])\subseteq(\FV(\varphi)\setminus\{y\})\cup\{y'\}
\subseteq\{x_1,\ldots,x_m,y'\}\), and \(\FV(y=y')=\{y,y'\}\); the free
variables of the matrix of \(\operatorname{Uni}_{\varphi}\) therefore lie in
\(\{x_1,\ldots,x_m,y,y'\}\), and its \(m+2\) quantifiers remove them.  In
\(\delta_f\), \(\FV(f(x_1,\ldots,x_m)=y)=\{x_1,\ldots,x_m,y\}\), the defined
biconditional has the free variables of its two sides, and the \(m+1\)
quantifiers remove them all.

Fix an \(\Lang\)-structure \(\mathcal M\) and an assignment \(s\).

\emph{Clause 1.}  Applying the clause for \(\forall\) once for each of
\(x_1,\ldots,x_m\) and the clause for \(\exists\) once,
\(\mathcal M,s\entails\operatorname{Tot}_{\varphi}\) iff for every
\(\bar a\in(\dom\mathcal M)^m\) there is \(b\) with
\(\mathcal M,s[\bar x\mapsto\bar a][y\mapsto b]\entails\varphi\).  By
\cref{thm:p02-w02b-p02-018-instance}(3) for the list \(x_1,\ldots,x_m,y\),
the last condition says \((\bar a,b)\in R^{\mathcal M}_{\varphi}\).

\emph{Clause 2.}  Applying the clause for \(\forall\) for each of
\(x_1,\ldots,x_m,y,y'\), \(\mathcal M,s\entails\operatorname{Uni}_{\varphi}\)
iff for all \(\bar a\), \(b\) and \(b'\), the assignment
\(u\defeq s[\bar x\mapsto\bar a][y\mapsto b][y'\mapsto b']\) satisfies
\((\varphi\land\varphi[y'/y])\to y=y'\).  The variables
\(x_1,\ldots,x_m,y,y'\) are pairwise distinct, so by Fact U2,
\(u(x_i)=a_i\), \(u(y)=b\) and \(u(y')=b'\).
\begin{enumerate}
 \item \(u\) and \(s[\bar x\mapsto\bar a][y\mapsto b]\) agree on
       \(\{x_1,\ldots,x_m,y\}\supseteq\FV(\varphi)\), so by
       \cref{lem:fol-coincidence} and
       \cref{thm:p02-w02b-p02-018-instance}(3),
       \(\mathcal M,u\entails\varphi\) iff
       \((\bar a,b)\in R^{\mathcal M}_{\varphi}\).
 \item By \cref{thm:fol-substitution}, since \(\sem{y'}^{\mathcal M}_u=b'\),
       \(\mathcal M,u\entails\varphi[y'/y]\) iff
       \(\mathcal M,u[y\mapsto b']\entails\varphi\).  The assignment
       \(u[y\mapsto b']\) sends each \(x_i\) to \(a_i\) and \(y\) to \(b'\), so
       it agrees with \(s[\bar x\mapsto\bar a][y\mapsto b']\) on
       \(\FV(\varphi)\); as in item 1, this holds iff
       \((\bar a,b')\in R^{\mathcal M}_{\varphi}\).
 \item \(\mathcal M,u\entails y=y'\) iff \(b=b'\).
\end{enumerate}
By the clauses for \(\to\) and \(\land\), \(u\) satisfies the matrix iff
\((\bar a,b)\in R^{\mathcal M}_{\varphi}\) and
\((\bar a,b')\in R^{\mathcal M}_{\varphi}\) together imply \(b=b'\).  Holding
for all \(\bar a,b,b'\), this says that each \(\bar a\) has at most one
\(b\) with \((\bar a,b)\in R^{\mathcal M}_{\varphi}\).

\emph{Clause 3.}  ``Exactly one'' is ``at least one'' and ``at most one'';
combine clauses 1 and 2 with
\cref{def:p02-w02b-p02-018-defined-function}.
\end{proof}
\end{annexitem}

\begin{annexitem}{annex:p02-w02b-p02-018-fn-expansion}{thm:p02-w02b-p02-018-fn-expansion}{The defining expansion}
\begin{proof}
\emph{Clause 1.}  Let \(\mathcal N\) be an \(\Lang_f\)-structure, put
\(\mathcal M\defeq\mathcal N\mathbin{\upharpoonright}\Lang\), and fix an
assignment \(s\).  Since \(\delta_f\) is a sentence
(\cref{thm:p02-w02b-p02-018-unique-value}), applying the clause for
\(\forall\) for each of \(x_1,\ldots,x_m,y\) gives:
\(\mathcal N\entails\delta_f\) iff for all \(\bar a\in(\dom\mathcal N)^m\)
and \(b\in\dom\mathcal N\), the assignment
\(s'\defeq s[\bar x\mapsto\bar a][y\mapsto b]\) satisfies both
\(f(x_1,\ldots,x_m)=y\) and \(\varphi\) or neither (the computation of the
defined biconditional in \cref{annex:p02-w02b-p02-018-consequences}).  By
Fact U2, \(s'(x_i)=a_i\) and \(s'(y)=b\).  By the clause for equality and the
definition of term value, \(\mathcal N,s'\entails f(x_1,\ldots,x_m)=y\) iff
\(f^{\mathcal N}(\bar a)=b\).  By \cref{lem:reduct-invariance} and
\cref{thm:p02-w02b-p02-018-instance}(3),
\(\mathcal N,s'\entails\varphi\) iff \(\mathcal M,s'\entails\varphi\) iff
\((\bar a,b)\in R^{\mathcal M}_{\varphi}\).  Hence
\(\mathcal N\entails\delta_f\) iff
\[
 \text{for all }\bar a\text{ and }b:\quad
 f^{\mathcal N}(\bar a)=b\ \Longleftrightarrow\
 (\bar a,b)\in R^{\mathcal M}_{\varphi}.
 \tag{\(\ast\)}
\]
Suppose \((\ast)\) and fix \(\bar a\).  Taking \(b=f^{\mathcal N}(\bar a)\)
gives \((\bar a,f^{\mathcal N}(\bar a))\in R^{\mathcal M}_{\varphi}\), and any
\(b\) with \((\bar a,b)\in R^{\mathcal M}_{\varphi}\) equals
\(f^{\mathcal N}(\bar a)\).  So \(\varphi\) is functional in \(\mathcal M\)
and \(f^{\mathcal M}_{\varphi}(\bar a)=f^{\mathcal N}(\bar a)\).  Conversely,
if \(\varphi\) is functional in \(\mathcal M\) and
\(f^{\mathcal N}=f^{\mathcal M}_{\varphi}\), then for all \(\bar a,b\):
\((\bar a,b)\in R^{\mathcal M}_{\varphi}\) iff
\(b=f^{\mathcal M}_{\varphi}(\bar a)=f^{\mathcal N}(\bar a)\), which is
\((\ast)\).  Finally, when \(\varphi\) is functional in \(\mathcal M\), the
structures \(\mathcal N\) and \(\mathcal M^{f}_{\varphi}\) are both
\(\Lang_f\)-expansions of \(\mathcal M\); they have the same domain and
agree on \(\Lang\), so they are equal iff
\(f^{\mathcal N}=f^{\mathcal M}_{\varphi}\).  For \(m=0\) the tuple
\(\bar a\) is empty and \(f^{\mathcal N}\) is identified with its value; the
argument is unchanged.

\emph{Clause 2.}  If \(\mathcal N\) is an \(\Lang_f\)-expansion of
\(\mathcal M\) satisfying \(\delta_f\), then
\(\mathcal N\mathbin{\upharpoonright}\Lang=\mathcal M\), so by clause 1
\(\varphi\) is functional in \(\mathcal M\) and
\(\mathcal N=\mathcal M^{f}_{\varphi}\); by
\cref{thm:p02-w02b-p02-018-unique-value}(3),
\(\mathcal M\entails\operatorname{Tot}_{\varphi}\) and
\(\mathcal M\entails\operatorname{Uni}_{\varphi}\).  This proves the
left-to-right direction and the uniqueness.  Conversely, if \(\mathcal M\)
satisfies both sentences, \(\varphi\) is functional in \(\mathcal M\) by
\cref{thm:p02-w02b-p02-018-unique-value}(3), so \(\mathcal M^{f}_{\varphi}\)
is defined; its reduct to \(\Lang\) is \(\mathcal M\) and its interpretation
of \(f\) is \(f^{\mathcal M}_{\varphi}\), so it satisfies \(\delta_f\) by
clause 1.

\emph{Clause 3.}  Let \(\mathcal M\) be an \(\Lang\)-model of \(T\).  The
premise gives \(\mathcal M\entails\operatorname{Tot}_{\varphi}\) and
\(\mathcal M\entails\operatorname{Uni}_{\varphi}\), and clause 2 gives
exactly one \(\Lang_f\)-expansion satisfying \(\delta_f\), namely
\(\mathcal M^{f}_{\varphi}\).  Each member of \(T\) is an \(\Lang\)-sentence
true in \(\mathcal M\), hence true in \(\mathcal M^{f}_{\varphi}\) by
\cref{lem:reduct-invariance}; so \(\mathcal M^{f}_{\varphi}\) is a model of
\(T\cup\{\delta_f\}\).

\emph{Clause 4.}  Let \(\mathcal N\) and \(s\) satisfy \(\delta_f\).  By
clause 1, \(\varphi\) is functional in
\(\mathcal N\mathbin{\upharpoonright}\Lang\); by
\cref{thm:p02-w02b-p02-018-unique-value}(3) that reduct satisfies
\(\operatorname{Tot}_{\varphi}\) and \(\operatorname{Uni}_{\varphi}\), and by
\cref{lem:reduct-invariance} so does \(\mathcal N\) under \(s\).
\end{proof}
\end{annexitem}

\begin{annexitem}{annex:p02-w02b-p02-018-graph-formula}{thm:p02-w02b-p02-018-graph-formula}{Graph formulas}
\begin{proof}
Both clauses are proved together by induction on the construction of \(t\),
for every \(v\notin\operatorname{Var}(t)\) and, in clause 2, every
assignment \(s\).  In clause 2 write
\(\mathcal N\defeq\mathcal M^{f}_{\varphi}\); recall
\(\mathcal N\mathbin{\upharpoonright}\Lang=\mathcal M\).

\emph{Case 1: \(f\) does not occur in \(t\).}  By induction on \(t\), a
term of \(\Lang_f\) in which \(f\) does not occur is an \(\Lang\)-term: a
variable is one, and \(g(t_1,\ldots,t_k)\) with \(g\ne f\) has \(g\in\Lang\)
and \(\Lang\)-terms as arguments.  So \(v=t\) is an atomic
\(\Lang\)-formula with \(\FV(v=t)=\{v\}\cup\operatorname{Var}(t)\).  For
clause 2, \(\mathcal M,s\entails v=t\) iff
\(s(v)=\sem{t}^{\mathcal M}_s\), and
\(\sem{t}^{\mathcal M}_s=\sem{t}^{\mathcal N}_s\) by
\cref{lem:reduct-invariance}.

\emph{Cases 2 and 3.}  Let \(t=h(t_1,\ldots,t_j)\), where either
\(h=g\in\Lang\), \(j=k\) and \(\beta\) is \(v=g(w_1,\ldots,w_j)\) (case 2),
or \(h=f\), \(j=m\) and \(\beta\) is \(\varphi(w_1,\ldots,w_j,v)\) (case 3).
The variables \(w_1,\ldots,w_j\) are pairwise distinct, lie outside
\(\operatorname{Var}(t)\supseteq\operatorname{Var}(t_i)\), and differ from
\(v\).

\emph{Clause 1.}  By the induction hypothesis,
\(\operatorname{Gr}_{t_i}(w_i)\in\mathsf{Fm}(\Lang)\) with free variables in
\(\operatorname{Var}(t_i)\cup\{w_i\}\).  In case 2, \(\beta\) is an atomic
\(\Lang\)-formula with \(\FV(\beta)=\{v,w_1,\ldots,w_j\}\).  In case 3,
\cref{thm:p02-w02b-p02-018-instance}(1), applied to the terms
\(w_1,\ldots,w_j,v\), gives \(\beta\in\mathsf{Fm}(\Lang)\) with
\(\FV(\beta)\subseteq\{w_1,\ldots,w_j,v\}\).  The conjunction and the
existential quantifiers keep the formula in \(\mathsf{Fm}(\Lang)\), and its
free variables lie in
\[
 \Bigl(\textstyle\bigcup_i\operatorname{Var}(t_i)\cup\{w_1,\ldots,w_j,v\}\Bigr)
 \setminus\{w_1,\ldots,w_j\}
 \subseteq\operatorname{Var}(t)\cup\{v\},
\]
because \(\operatorname{Var}(t)=\bigcup_i\operatorname{Var}(t_i)\).

\emph{Clause 2.}  By Fact U3, \(\mathcal M,s\entails\operatorname{Gr}_{t}(v)\)
iff there are \(c_1,\ldots,c_j\) such that
\(s_c\defeq s[w_1\mapsto c_1]\cdots[w_j\mapsto c_j]\) satisfies every
\(\operatorname{Gr}_{t_i}(w_i)\) and \(\beta\).  By Fact U2,
\(s_c(w_i)=c_i\), and \(s_c\) agrees with \(s\) at every other variable, in
particular at \(v\) and on \(\operatorname{Var}(t)\).  Since
\(w_i\notin\operatorname{Var}(t_i)\), the induction hypothesis at \(s_c\)
gives \(\mathcal M,s_c\entails\operatorname{Gr}_{t_i}(w_i)\) iff
\(c_i=\sem{t_i}^{\mathcal N}_{s_c}\), and term coincidence gives
\(\sem{t_i}^{\mathcal N}_{s_c}=\sem{t_i}^{\mathcal N}_{s}\); write
\(d_i\) for this value.  So
the conditions on the \(\operatorname{Gr}_{t_i}(w_i)\) hold exactly when
\(c_i=d_i\) for every \(i\), and
\(\mathcal M,s\entails\operatorname{Gr}_{t}(v)\) iff
\(\mathcal M,s_d\entails\beta\), where
\(s_d\defeq s[w_1\mapsto d_1]\cdots[w_j\mapsto d_j]\).

In case 2, \(\mathcal M,s_d\entails v=g(w_1,\ldots,w_j)\) iff
\(s(v)=g^{\mathcal M}(d_1,\ldots,d_j)\).  Since \(\mathcal N\) is an expansion
of \(\mathcal M\), \(g^{\mathcal N}=g^{\mathcal M}\), and
\(g^{\mathcal N}(d_1,\ldots,d_j)=\sem{t}^{\mathcal N}_s\).

In case 3, \cref{thm:p02-w02b-p02-018-instance}(2), applied in
\(\mathcal M\) at \(s_d\) to the terms \(w_1,\ldots,w_j,v\), whose values
are \(d_1,\ldots,d_j,s(v)\), gives \(\mathcal M,s_d\entails\beta\) iff
\(\mathcal M,s_d[\bar x\mapsto\bar d][y\mapsto s(v)]\entails\varphi\), and
\cref{thm:p02-w02b-p02-018-instance}(3) turns this into
\((\bar d,s(v))\in R^{\mathcal M}_{\varphi}\).  Since \(\varphi\) is
functional in \(\mathcal M\), this holds iff
\(s(v)=f^{\mathcal M}_{\varphi}(\bar d)=f^{\mathcal N}(\bar d)
=\sem{t}^{\mathcal N}_s\).  For \(m=0\) there are no \(w_i\), \(s_d=s\), and
the same steps give \(s(v)=f^{\mathcal N}=\sem{f}^{\mathcal N}_s\).
\end{proof}
\end{annexitem}

\begin{annexitem}{annex:p02-w02b-p02-018-fn-syntax}{thm:p02-w02b-p02-018-fn-syntax}{The function translation lands in the old signature}
\begin{proof}
By strong induction on \(c(\psi)\), following the clauses of
\cref{def:p02-w02b-p02-018-function-translation}.

\emph{Atomic, without \(f\).}  The terms of \(\psi\) are
\(\Lang\)-terms, as shown in case 1 of
\cref{annex:p02-w02b-p02-018-graph-formula}, and the two signatures have
the same relation symbols, so
\(\psi^{\sharp}=\psi\in\mathsf{Fm}(\Lang)\).

\emph{Atomic, with \(f\).}  For \(r=t\): \(v_1\notin\operatorname{Var}(r)\)
and \(v_2\notin\operatorname{Var}(t)\), so both graph formulas are defined;
by \cref{thm:p02-w02b-p02-018-graph-formula}(1) they are \(\Lang\)-formulas
with free variables in \(\operatorname{Var}(r)\cup\{v_1\}\) and
\(\operatorname{Var}(t)\cup\{v_2\}\).  The atom \(v_1=v_2\) is an
\(\Lang\)-formula, and the two quantifiers remove \(v_1\) and \(v_2\), so
\(\FV((r=t)^{\sharp})\subseteq\operatorname{Var}(r)\cup\operatorname{Var}(t)
=\FV(r=t)\).

For \(\psi=S(t_1,\ldots,t_k)\) with \(S\in\Lang\) and \(f\) occurring in
\(\psi\): the variables \(v_1,\ldots,v_k\) are pairwise distinct and lie
outside \(\operatorname{Var}(t_1)\cup\cdots\cup\operatorname{Var}(t_k)\), so
\(v_i\notin\operatorname{Var}(t_i)\) and each \(\operatorname{Gr}_{t_i}(v_i)\)
is defined.  By \cref{thm:p02-w02b-p02-018-graph-formula}(1),
\(\operatorname{Gr}_{t_i}(v_i)\in\mathsf{Fm}(\Lang)\) with
\(\FV(\operatorname{Gr}_{t_i}(v_i))\subseteq\operatorname{Var}(t_i)\cup\{v_i\}\).
Since \(S\) is a relation symbol of \(\Lang\) of arity \(k\) and the
variables \(v_1,\ldots,v_k\) are \(\Lang\)-terms, \(S(v_1,\ldots,v_k)\) is an
atomic \(\Lang\)-formula with \(\FV(S(v_1,\ldots,v_k))=\{v_1,\ldots,v_k\}\).
The \(k\) conjunctions, formed right-associated, and the \(k\) existential
quantifiers keep the formula in \(\mathsf{Fm}(\Lang)\).  The free variables
of the conjunction lie in
\(\operatorname{Var}(t_1)\cup\cdots\cup\operatorname{Var}(t_k)
\cup\{v_1,\ldots,v_k\}\), and the quantifiers
\(\exists v_1\cdots\exists v_k\) remove \(v_1,\ldots,v_k\), so
\[
 \FV\bigl((S(t_1,\ldots,t_k))^{\sharp}\bigr)
 \subseteq\operatorname{Var}(t_1)\cup\cdots\cup\operatorname{Var}(t_k)
 =\FV(S(t_1,\ldots,t_k)).
\]

\emph{Falsum.}  \(\bottom^{\sharp}=\bottom\in\mathsf{Fm}(\Lang)\) and
\(\FV(\bottom)=\varnothing\).

\emph{Negation.}  For \(\psi=\neg\chi\), \(c(\chi)<c(\psi)\), and the
induction hypothesis gives \(\chi^{\sharp}\in\mathsf{Fm}(\Lang)\) with
\(\FV(\chi^{\sharp})\subseteq\FV(\chi)\).  Hence
\(\psi^{\sharp}=\neg\chi^{\sharp}\in\mathsf{Fm}(\Lang)\) and
\(\FV(\neg\chi^{\sharp})=\FV(\chi^{\sharp})\subseteq\FV(\chi)=\FV(\psi)\).

\emph{Binary connectives.}  For \(\psi=(\chi\circ\eta)\) with
\(\circ\in\{\land,\lor,\to\}\), \(c(\chi),c(\eta)<c(\psi)\), and the
induction hypothesis gives \(\chi^{\sharp},\eta^{\sharp}\in\mathsf{Fm}(\Lang)\)
with \(\FV(\chi^{\sharp})\subseteq\FV(\chi)\) and
\(\FV(\eta^{\sharp})\subseteq\FV(\eta)\).  Hence
\(\psi^{\sharp}=(\chi^{\sharp}\circ\eta^{\sharp})\in\mathsf{Fm}(\Lang)\) and
\(\FV(\psi^{\sharp})=\FV(\chi^{\sharp})\cup\FV(\eta^{\sharp})
\subseteq\FV(\chi)\cup\FV(\eta)=\FV(\psi)\).

\emph{Quantifiers.}  For \(\psi=Qu\,\chi\) with \(Q\in\{\forall,\exists\}\),
\(c(\chi)<c(\psi)\), and the induction hypothesis gives
\(\chi^{\sharp}\in\mathsf{Fm}(\Lang)\) with
\(\FV(\chi^{\sharp})\subseteq\FV(\chi)\).  Hence
\(\psi^{\sharp}=Qu\,\chi^{\sharp}\in\mathsf{Fm}(\Lang)\) and
\(\FV(Qu\,\chi^{\sharp})=\FV(\chi^{\sharp})\setminus\{u\}
\subseteq\FV(\chi)\setminus\{u\}=\FV(\psi)\).

\emph{Old formulas.}  If \(\psi\in\mathsf{Fm}(\Lang)\), then \(f\) occurs in
no atomic subformula of \(\psi\), so only clause 1 applies to atoms, and
the remaining clauses rebuild \(\psi\) from translations of immediate
subformulas, which are fixed by the induction hypothesis.
\end{proof}
\end{annexitem}

\begin{annexitem}{annex:p02-w02b-p02-018-fn-satisfaction}{thm:p02-w02b-p02-018-fn-satisfaction}{Function satisfaction translation}
\begin{proof}
Let \(\varphi\) be functional in \(\mathcal M\) and write
\(\mathcal N\defeq\mathcal M^{f}_{\varphi}\).  We prove the equivalence by
strong induction on \(c(\psi)\), for every assignment \(s\).

\emph{Atomic, without \(f\).}  Then \(\psi\in\mathsf{Fm}(\Lang)\) and
\(\psi^{\sharp}=\psi\); \cref{lem:reduct-invariance} gives
\(\mathcal N,s\entails\psi\) iff \(\mathcal M,s\entails\psi\).

\emph{The atom \(r=t\) with \(f\).}  By Fact U3,
\(\mathcal M,s\entails(r=t)^{\sharp}\) iff there are \(c_1,c_2\) such that
\(s_c\defeq s[v_1\mapsto c_1][v_2\mapsto c_2]\) satisfies
\(\operatorname{Gr}_{r}(v_1)\), \(\operatorname{Gr}_{t}(v_2)\) and
\(v_1=v_2\).  By Fact U2, \(s_c(v_1)=c_1\) and \(s_c(v_2)=c_2\).  By
\cref{thm:p02-w02b-p02-018-graph-formula}(2) at \(s_c\) and term coincidence
(\(v_1,v_2\notin\operatorname{Var}(r)\cup\operatorname{Var}(t)\)), the first
two conditions say \(c_1=\sem{r}^{\mathcal N}_s\) and
\(c_2=\sem{t}^{\mathcal N}_s\); the third says \(c_1=c_2\).  If witnesses
exist, then \(\sem{r}^{\mathcal N}_s=c_1=c_2=\sem{t}^{\mathcal N}_s\);
conversely, if \(\sem{r}^{\mathcal N}_s=\sem{t}^{\mathcal N}_s\), the
choice \(c_1\defeq\sem{r}^{\mathcal N}_s\), \(c_2\defeq\sem{t}^{\mathcal N}_s\)
meets all three conditions.  So
\(\mathcal M,s\entails(r=t)^{\sharp}\) iff
\(\sem{r}^{\mathcal N}_s=\sem{t}^{\mathcal N}_s\), that is, iff
\(\mathcal N,s\entails r=t\).

\emph{The atom \(S(t_1,\ldots,t_k)\) with \(f\).}  Let
\(\psi=S(t_1,\ldots,t_k)\) with \(S\in\Lang\) of arity \(k\) and \(f\)
occurring in \(\psi\), and let \(v_1,\ldots,v_k\) be the variables of
\cref{def:p02-w02b-p02-018-function-translation}: pairwise distinct and
outside \(\operatorname{Var}(t_1)\cup\cdots\cup\operatorname{Var}(t_k)\).
For each \(i\) put \(d_i\defeq\sem{t_i}^{\mathcal N}_s\).  By Fact U3 with
\(j=k\), \(u_i=v_i\), \(\alpha_i=\operatorname{Gr}_{t_i}(v_i)\) and
\(\beta=S(v_1,\ldots,v_k)\), \(\mathcal M,s\entails\psi^{\sharp}\) iff there
are \(c_1,\ldots,c_k\in\dom\mathcal M\) such that the assignment
\(s_c\defeq s[v_1\mapsto c_1]\cdots[v_k\mapsto c_k]\) satisfies
\(\operatorname{Gr}_{t_i}(v_i)\) for every \(i\) and satisfies
\(S(v_1,\ldots,v_k)\).  Fix such \(c_1,\ldots,c_k\) for the moment and
examine the conditions one at a time.
\begin{enumerate}
 \item By Fact U2, \(s_c(v_i)=c_i\) for every \(i\), and \(s_c\) agrees with
       \(s\) at every variable other than \(v_1,\ldots,v_k\); in particular
       \(s_c\) and \(s\) agree on each \(\operatorname{Var}(t_i)\).
 \item Since \(v_i\notin\operatorname{Var}(t_i)\) and \(\varphi\) is
       functional in \(\mathcal M\),
       \cref{thm:p02-w02b-p02-018-graph-formula}(2) at the assignment
       \(s_c\) gives \(\mathcal M,s_c\entails\operatorname{Gr}_{t_i}(v_i)\)
       iff \(s_c(v_i)=\sem{t_i}^{\mathcal N}_{s_c}\), that is, iff
       \(c_i=\sem{t_i}^{\mathcal N}_{s_c}\).  By item 1 and term coincidence
       in \(\mathcal N\), \(\sem{t_i}^{\mathcal N}_{s_c}=\sem{t_i}^{\mathcal N}_{s}=d_i\).
       So this condition is \(c_i=d_i\).
 \item By the atomic clause in \(\mathcal M\) and item 1,
       \(\mathcal M,s_c\entails S(v_1,\ldots,v_k)\) iff
       \((s_c(v_1),\ldots,s_c(v_k))=(c_1,\ldots,c_k)\in S^{\mathcal M}\).
\end{enumerate}
If witnesses \(c_1,\ldots,c_k\) exist, items 2 and 3 give \(c_i=d_i\) for
every \(i\) and \((c_1,\ldots,c_k)\in S^{\mathcal M}\), so
\((d_1,\ldots,d_k)\in S^{\mathcal M}\).  Conversely, if
\((d_1,\ldots,d_k)\in S^{\mathcal M}\), the choice \(c_i\defeq d_i\) meets
every condition in items 2 and 3.  Hence
\(\mathcal M,s\entails\psi^{\sharp}\) iff
\((d_1,\ldots,d_k)\in S^{\mathcal M}\).  Since \(\mathcal N\) is an
\(\Lang_f\)-expansion of \(\mathcal M\) and \(S\in\Lang\),
\(S^{\mathcal N}=S^{\mathcal M}\), and by the atomic clause in \(\mathcal N\),
\((d_1,\ldots,d_k)=(\sem{t_1}^{\mathcal N}_s,\ldots,\sem{t_k}^{\mathcal N}_s)
\in S^{\mathcal N}\) iff \(\mathcal N,s\entails S(t_1,\ldots,t_k)\).

\emph{Falsum.}  \(\bottom^{\sharp}=\bottom\), and neither \(\mathcal N\)
nor \(\mathcal M\) satisfies \(\bottom\) under any assignment.

In the remaining cases the immediate subformulas \(\chi\) (and \(\eta\))
have lower complexity than \(\psi\), so the induction hypothesis applies to
them at every assignment.

\emph{Negation.}  \(\mathcal N,s\entails\neg\chi\) iff
\(\mathcal N,s\not\entails\chi\) iff, by the induction hypothesis at \(s\),
\(\mathcal M,s\not\entails\chi^{\sharp}\) iff
\(\mathcal M,s\entails\neg\chi^{\sharp}\), and
\(\neg\chi^{\sharp}=(\neg\chi)^{\sharp}\).

\emph{Binary connectives.}  The induction hypothesis at \(s\) gives
\(\mathcal N,s\entails\chi\) iff \(\mathcal M,s\entails\chi^{\sharp}\), and
\(\mathcal N,s\entails\eta\) iff \(\mathcal M,s\entails\eta^{\sharp}\).
\begin{enumerate}
 \item \(\mathcal N,s\entails(\chi\land\eta)\) iff
       \(\mathcal N,s\entails\chi\) and \(\mathcal N,s\entails\eta\); iff
       \(\mathcal M,s\entails\chi^{\sharp}\) and
       \(\mathcal M,s\entails\eta^{\sharp}\); iff
       \(\mathcal M,s\entails(\chi^{\sharp}\land\eta^{\sharp})
       =(\chi\land\eta)^{\sharp}\).
 \item \(\mathcal N,s\entails(\chi\lor\eta)\) iff
       \(\mathcal N,s\entails\chi\) or \(\mathcal N,s\entails\eta\); iff
       \(\mathcal M,s\entails\chi^{\sharp}\) or
       \(\mathcal M,s\entails\eta^{\sharp}\); iff
       \(\mathcal M,s\entails(\chi^{\sharp}\lor\eta^{\sharp})
       =(\chi\lor\eta)^{\sharp}\).
 \item \(\mathcal N,s\entails(\chi\to\eta)\) iff
       \(\mathcal N,s\not\entails\chi\) or \(\mathcal N,s\entails\eta\); iff
       \(\mathcal M,s\not\entails\chi^{\sharp}\) or
       \(\mathcal M,s\entails\eta^{\sharp}\); iff
       \(\mathcal M,s\entails(\chi^{\sharp}\to\eta^{\sharp})
       =(\chi\to\eta)^{\sharp}\).
\end{enumerate}

\emph{Quantifiers.}  The structure \(\mathcal N\) is an expansion of
\(\mathcal M\), so \(\dom\mathcal N=\dom\mathcal M\).  Let
\(\psi=\forall u\,\chi\).  By the clause for \(\forall\),
\(\mathcal N,s\entails\psi\) iff \(\mathcal N,s[u\mapsto a]\entails\chi\)
for every \(a\in\dom\mathcal M\).  The induction hypothesis at each
assignment \(s[u\mapsto a]\) turns this into:
\(\mathcal M,s[u\mapsto a]\entails\chi^{\sharp}\) for every
\(a\in\dom\mathcal M\), which is
\(\mathcal M,s\entails\forall u\,\chi^{\sharp}\), and
\(\forall u\,\chi^{\sharp}=\psi^{\sharp}\).  Let \(\psi=\exists u\,\chi\).
By the clause for \(\exists\), \(\mathcal N,s\entails\psi\) iff there is
\(a\in\dom\mathcal M\) with \(\mathcal N,s[u\mapsto a]\entails\chi\); by the
induction hypothesis at \(s[u\mapsto a]\), iff there is such an \(a\) with
\(\mathcal M,s[u\mapsto a]\entails\chi^{\sharp}\); and this is
\(\mathcal M,s\entails\exists u\,\chi^{\sharp}\), with
\(\exists u\,\chi^{\sharp}=\psi^{\sharp}\).

\emph{The model-of-\(T\) clause.}  If \(\mathcal M\) is a model of \(T\) and
the premise holds, then \(\mathcal M\) satisfies both unique-value
sentences, so \(\varphi\) is functional in \(\mathcal M\) by
\cref{thm:p02-w02b-p02-018-unique-value}(3).
\end{proof}
\end{annexitem}

\begin{annexitem}{annex:p02-w02b-p02-018-fn-conservativity}{thm:p02-w02b-p02-018-fn-conservativity}{Conservativity for a function symbol}
\begin{proof}
\emph{Clause 1.}  Let \(\mathcal N\) and \(s\) satisfy \(\delta_f\), and put
\(\mathcal M\defeq\mathcal N\mathbin{\upharpoonright}\Lang\).  By
\cref{thm:p02-w02b-p02-018-fn-expansion}(1), \(\varphi\) is functional
in \(\mathcal M\) and \(\mathcal N=\mathcal M^{f}_{\varphi}\).  By
\cref{thm:p02-w02b-p02-018-fn-satisfaction},
\(\mathcal N,s\entails\psi\) iff \(\mathcal M,s\entails\psi^{\sharp}\), and
since \(\psi^{\sharp}\in\mathsf{Fm}(\Lang)\)
(\cref{thm:p02-w02b-p02-018-fn-syntax}),
\cref{lem:reduct-invariance} gives \(\mathcal M,s\entails\psi^{\sharp}\) iff
\(\mathcal N,s\entails\psi^{\sharp}\).  So \(\mathcal N,s\) satisfies both
or neither of \(\psi\) and \(\psi^{\sharp}\), which is the biconditional.

\emph{Clause 2, left to right.}  Suppose
\(T\cup\{\delta_f\}\cup\Gamma\entails\psi\), and let the
\(\Lang\)-structure \(\mathcal M\) and assignment \(s\) satisfy every member
of \(T\cup\Gamma^{\sharp}\).  Then \(\mathcal M\) is a model of \(T\), so by
the premise and \cref{thm:p02-w02b-p02-018-unique-value}(3) \(\varphi\) is
functional in \(\mathcal M\).  By
\cref{thm:p02-w02b-p02-018-fn-expansion}(3),
\(\mathcal M^{f}_{\varphi}\) satisfies \(T\cup\{\delta_f\}\); by
\cref{thm:p02-w02b-p02-018-fn-satisfaction}, it satisfies every
\(\gamma\in\Gamma\) under \(s\), because \(\mathcal M,s\entails\gamma^{\sharp}\).
Hence \(\mathcal M^{f}_{\varphi},s\entails\psi\), and the same theorem gives
\(\mathcal M,s\entails\psi^{\sharp}\).  As \(\mathcal M\) and \(s\) were
arbitrary, \(T\cup\Gamma^{\sharp}\entails\psi^{\sharp}\).

\emph{Clause 2, right to left.}  Suppose
\(T\cup\Gamma^{\sharp}\entails\psi^{\sharp}\), and let the
\(\Lang_f\)-structure \(\mathcal N\) and assignment \(s\) satisfy every
member of \(T\cup\{\delta_f\}\cup\Gamma\).  Put
\(\mathcal M\defeq\mathcal N\mathbin{\upharpoonright}\Lang\).  By
\cref{thm:p02-w02b-p02-018-fn-expansion}(1), \(\varphi\) is functional
in \(\mathcal M\) and \(\mathcal N=\mathcal M^{f}_{\varphi}\).  By
\cref{lem:reduct-invariance}, \(\mathcal M,s\) satisfies \(T\); by
\cref{thm:p02-w02b-p02-018-fn-satisfaction}, it satisfies
\(\Gamma^{\sharp}\).  Hence \(\mathcal M,s\entails\psi^{\sharp}\), and the
same theorem gives \(\mathcal N,s\entails\psi\).  As \(\mathcal N\) and
\(s\) were arbitrary, \(T\cup\{\delta_f\}\cup\Gamma\entails\psi\).  This
direction does not use the premise.

For \(\psi\in\mathsf{Fm}(\Lang)\) and \(\Gamma=\varnothing\),
\(\psi^{\sharp}=\psi\) by
\cref{thm:p02-w02b-p02-018-fn-syntax}, which gives the
displayed special case.

\emph{Clause 3.}  Assume the Axiom of Choice.  Every structure and
assignment satisfying \(T\cup\{\delta_f\}\) satisfy \(\delta_f\), so clause 1
gives \(T\cup\{\delta_f\}\entails\psi\leftrightarrow\psi^{\sharp}\) in the
signature \(\Lang_f\), and \cref{thm:fol-completeness} in \(\Lang_f\), applied
to the formula \(((\psi\to\psi^{\sharp})\land(\psi^{\sharp}\to\psi))\), gives
the derivation.  If
\(T\cup\{\delta_f\}\cup\Gamma\derives_{\mathsf{ND}_{\mathsf{FOL}}}\psi\), then
\cref{thm:fol-soundness} gives \(T\cup\{\delta_f\}\cup\Gamma\entails\psi\),
clause 2 gives \(T\cup\Gamma^{\sharp}\entails\psi^{\sharp}\), and
\cref{thm:fol-completeness} in \(\Lang\) gives
\(T\cup\Gamma^{\sharp}\derives_{\mathsf{ND}_{\mathsf{FOL}}}\psi^{\sharp}\).
Conversely, if
\(T\cup\Gamma^{\sharp}\derives_{\mathsf{ND}_{\mathsf{FOL}}}\psi^{\sharp}\) in
\(\Lang\), then \cref{thm:fol-soundness} gives
\(T\cup\Gamma^{\sharp}\entails\psi^{\sharp}\), clause 2 gives
\(T\cup\{\delta_f\}\cup\Gamma\entails\psi\), and \cref{thm:fol-completeness}
in \(\Lang_f\) gives
\(T\cup\{\delta_f\}\cup\Gamma\derives_{\mathsf{ND}_{\mathsf{FOL}}}\psi\).

\emph{Clause 4.}  If \(T\) entails both unique-value sentences, clause 2
gives conservativity for \(\Lang\)-sentences.  Conversely, suppose that for
every \(\sigma\in\mathsf{Sent}(\Lang)\),
\(T\cup\{\delta_f\}\entails\sigma\) implies \(T\entails\sigma\).  By
\cref{thm:p02-w02b-p02-018-fn-expansion}(4), every structure and
assignment satisfying \(T\cup\{\delta_f\}\) satisfy
\(\operatorname{Tot}_{\varphi}\) and \(\operatorname{Uni}_{\varphi}\); that
is, \(T\cup\{\delta_f\}\entails\operatorname{Tot}_{\varphi}\) and
\(T\cup\{\delta_f\}\entails\operatorname{Uni}_{\varphi}\).  Both are
\(\Lang\)-sentences by \cref{thm:p02-w02b-p02-018-unique-value}, so the
assumed conservativity, applied to \(\sigma=\operatorname{Tot}_{\varphi}\)
and to \(\sigma=\operatorname{Uni}_{\varphi}\), gives
\(T\entails\operatorname{Tot}_{\varphi}\) and
\(T\entails\operatorname{Uni}_{\varphi}\).
\end{proof}
\end{annexitem}

\begin{annexitem}{annex:p02-w02b-p02-018-iterated}{thm:p02-w02b-p02-018-iterated}{Finite sequences of definitions}
\begin{proof}
For each \(j\), the results of this chapter for a single relation symbol,
or for a single function or constant symbol, hold with \(\Lang_{j-1}\) and
\(T_{j-1}\) in place of \(\Lang\) and \(T\), since their proofs use no
property of \(\Lang\) or \(T\) beyond the standing data.  Write
\(\tau_j:\mathsf{Fm}(\Lang_j)\to\mathsf{Fm}(\Lang_{j-1})\) for the translation
of step \(j\): \(\psi\mapsto\psi^{*}\) of
\cref{def:p02-w02b-p02-018-translation} if \(e_j\) is a relation symbol,
\(\psi\mapsto\psi^{\sharp}\) of
\cref{def:p02-w02b-p02-018-function-translation} otherwise, each formed from
the defining formula of step \(j\).  In both cases,
\(\{\delta_j\}\entails\psi\leftrightarrow\tau_j(\psi)\) for every
\(\psi\in\mathsf{Fm}(\Lang_j)\)
(\cref{thm:p02-w02b-p02-018-consequences}(2),
\cref{thm:p02-w02b-p02-018-fn-conservativity}(1)),
\(\FV(\tau_j(\psi))\subseteq\FV(\psi)\), and \(\tau_j(\psi)=\psi\) for
\(\psi\in\mathsf{Fm}(\Lang_{j-1})\)
(\cref{thm:p02-w02b-p02-018-translation-syntax},
\cref{thm:p02-w02b-p02-018-fn-syntax}).  We argue by
induction on \(k\).  For \(k=0\) all three clauses hold with
\(\psi^{[0]}\defeq\psi\).  Let \(k\geq1\) and assume the clauses for the
first \(k-1\) steps.

\emph{Clause 1.}  Let \(\mathcal M\) be an \(\Lang\)-model of \(T\).  By the
induction hypothesis it has exactly one \(\Lang_{k-1}\)-expansion
\(\mathcal M_{k-1}\) that is a model of \(T_{k-1}\).  If \(e_k\) is a
relation symbol, \cref{thm:p02-w02b-p02-018-consequences}(1) gives exactly
one \(\Lang_k\)-expansion of \(\mathcal M_{k-1}\) satisfying \(\delta_k\).  If
\(e_k\) is a function or constant symbol, \(\mathcal M_{k-1}\) is a model of
\(T_{k-1}\) and the premise holds relative to \(T_{k-1}\), so
\cref{thm:p02-w02b-p02-018-fn-expansion}(3) gives the same
conclusion.  That expansion satisfies \(T_{k-1}\) by
\cref{lem:reduct-invariance}, hence is a model of \(T_k\), and it is an
expansion of \(\mathcal M\).  For uniqueness, let \(\mathcal N\) be any
\(\Lang_k\)-expansion of \(\mathcal M\) that is a model of \(T_k\).  Its
reduct \(\mathcal N\mathbin{\upharpoonright}\Lang_{k-1}\) has the domain and
the \(\Lang\)-interpretations of \(\mathcal M\), so it is an
\(\Lang_{k-1}\)-expansion of \(\mathcal M\), and it is a model of
\(T_{k-1}\) by \cref{lem:reduct-invariance}; so it equals
\(\mathcal M_{k-1}\).  Then \(\mathcal N\) is an \(\Lang_k\)-expansion of
\(\mathcal M_{k-1}\) satisfying \(\delta_k\), hence the one found above.

\emph{Clause 2.}  Let \(\Gamma\cup\{\psi\}\subseteq\mathsf{Fm}(\Lang)\), a
subset of \(\mathsf{Fm}(\Lang_{k-1})\).  Since \(T_k=T_{k-1}\cup\{\delta_k\}\),
\cref{thm:p02-w02b-p02-018-consequences}(3) or
\cref{thm:p02-w02b-p02-018-fn-conservativity}(2), applied to step
\(k\), gives \(T_k\cup\Gamma\entails\psi\) iff
\(T_{k-1}\cup\tau_k[\Gamma]\entails\tau_k(\psi)\), where
\(\tau_k[\Gamma]=\{\tau_k(\gamma):\gamma\in\Gamma\}\).  Here
\(\tau_k[\Gamma]=\Gamma\) and \(\tau_k(\psi)=\psi\), so
\(T_k\cup\Gamma\entails\psi\) iff \(T_{k-1}\cup\Gamma\entails\psi\), and the
induction hypothesis gives iff \(T\cup\Gamma\entails\psi\).

\emph{Clause 3.}  For \(\psi\in\mathsf{Fm}(\Lang_k)\) put
\(\psi^{[k]}\defeq(\tau_k(\psi))^{[k-1]}\), where \((\cdot)^{[k-1]}\) is
given by the induction hypothesis for \(\tau_k(\psi)\in\mathsf{Fm}(\Lang_{k-1})\).
Then \(\psi^{[k]}\in\mathsf{Fm}(\Lang)\) and
\(\FV(\psi^{[k]})\subseteq\FV(\tau_k(\psi))\subseteq\FV(\psi)\).  Let the
\(\Lang_k\)-structure \(\mathcal N\) and assignment \(s\) satisfy \(T_k\).
They satisfy \(\delta_k\), so \(\mathcal N,s\entails\psi\) iff
\(\mathcal N,s\entails\tau_k(\psi)\).  The reduct
\(\mathcal N'\defeq\mathcal N\mathbin{\upharpoonright}\Lang_{k-1}\)
satisfies \(T_{k-1}\) under \(s\) by \cref{lem:reduct-invariance}, so by the
induction hypothesis \(\mathcal N',s\entails\tau_k(\psi)\) iff
\(\mathcal N',s\entails\psi^{[k]}\).  Both formulas belong to
\(\mathsf{Fm}(\Lang_{k-1})\), so \cref{lem:reduct-invariance} carries this
equivalence to \(\mathcal N\).  Chaining the two equivalences,
\(\mathcal N,s\) satisfies both or neither of \(\psi\) and \(\psi^{[k]}\);
hence \(T_k\entails\psi\leftrightarrow\psi^{[k]}\).
\end{proof}
\end{annexitem}
